% preambulo.tex - Paquetes y configuraciones comunes para ambos documentos
\documentclass[12pt,a4paper,oneside]{report}

\usepackage{ucaldas}  % Nuestra plantilla personalizada
\usepackage{etoolbox}

% Bibliografía
\usepackage[backend=biber,style=ieee,sorting=none]{biblatex}
\addbibresource{bibliography.bib}

% Gráficos
\graphicspath{{images/}}

% Tablas compactas y ajustadas al ancho útil de A4.
\AtBeginEnvironment{tabular}{\small\sloppy\setlength{\tabcolsep}{3pt}}
\AtBeginEnvironment{longtable}{\small\sloppy\setlength{\tabcolsep}{3pt}\setlength{\LTleft}{0pt}\setlength{\LTright}{0pt}}

% Configuración de minted para Python
\setminted{
    fontsize=\small,
    breaklines=true,
    breakafter=,
    frame=lines,
    framesep=4pt,
    linenos=true,
    numbersep=8pt,
    tabsize=4,
    bgcolor=LightBlue
}

% Configuración de pgfgantt
\usepackage{pgfgantt}
\pgfplotsset{compat=1.18}

% Configuración pgf-pie
\usetikzlibrary{pie}

% Comandos personalizados para el proyecto
\newcommand{\tacticalgrid}{\textsc{TacticalGrid}}
\newcommand{\rf}[1]{RF-#1}
\newcommand{\rfnf}[1]{RNF-#1}
\newcommand{\sectionref}[1]{Sección~\ref{#1}}
\newcommand{\figureref}[1]{Figura~\ref{#1}}
\newcommand{\tableref}[1]{Tabla~\ref{#1}}
\newcommand{\listingref}[1]{Listing~\ref{#1}}

% Abreviaturas
\newcommand{\ie}{i.\,e.}
\newcommand{\eg}{e.\,g.}
\newcommand{\cf}{c.\,f.}
\newcommand{\etal}{\textit{et al.}}

% Glosario técnico
\newcommand{\g}[1]{\textit{#1}}
\newcommand{\code}[1]{\texttt{#1}}
\newcommand{\var}[1]{\textit{#1}}
\newcommand{\diagramfigure}[3]{%
    \begin{figure}[H]
    \centering
    \IfFileExists{images/#1.pdf}{%
        \includegraphics[width=#2\textwidth,height=0.72\textheight,keepaspectratio]{images/#1.pdf}%
    }{%
        \fbox{\parbox[c][4cm][c]{0.82\textwidth}{\centering Diagrama no disponible: \texttt{#1.pdf}}}%
    }%
    \caption{#3}%
    \end{figure}%
}

\newcommand{\screenshotfigure}[3]{%
    \begin{figure}[H]
    \centering
    \IfFileExists{images/screenshots/#1.png}{%
        \includegraphics[width=#2\textwidth]{images/screenshots/#1.png}%
    }{\IfFileExists{images/screenshots/#1.jpg}{%
        \includegraphics[width=#2\textwidth]{images/screenshots/#1.jpg}%
    }{%
        \fbox{\parbox[c][6cm][c]{0.82\textwidth}{\centering Captura pendiente: \texttt{#1.png / #1.jpg}\\[0.2cm]#3}}%
    }}%
    \caption{#3}%
    \end{figure}%
}

% Referencias a RF/RNF
\newcommand{\RF}[1]{RF-#1}
\newcommand{\RNF}[1]{RNF-#1}

% Comandos para algoritmos
\newcommand{\astar}{A\textsuperscript{*}}
\newcommand{\alphabeta}{$\alpha$-$\beta$}
\newcommand{\minimax}{Minimax}
\newcommand{\bigO}[1]{\mathcal{O}(#1)}

% Referencias a figuras/tablas
\newcommand{\Fig}[1]{Figura~\ref{#1}}
\newcommand{\Tab}[1]{Tabla~\ref{#1}}
\newcommand{\Lst}[1]{Listing~\ref{#1}}
\newcommand{\Eq}[1]{Ecuación~\ref{#1}}
\newcommand{\Sec}[1]{Sección~\ref{#1}}
\newcommand{\App}[1]{Apéndice~\ref{#1}}

% Colores para código
\definecolor{codebg}{HTML}{1A1B2E}
\definecolor{codefg}{HTML}{C4C4E0}
\definecolor{keyword}{HTML}{6C63FF}
\newcommand{\kw}[1]{\textcolor{keyword}{\textbf{#1}}}

% Para tablas anchas
\newcommand{\rotatetable}[1]{\begin{sidewaystable}#1\end{sidewaystable}}
\usepackage{rotating}